
\section{量子力学：一维谐振子在 $\delta$ 脉冲下的跃迁——相干态与 Poisson 分布}
\label{sec:delta-pulse-ho}

\begin{modelbox}[一维谐振子在 $\delta$ 脉冲下的跃迁]
谐振子基态受瞬时脉冲后的跃迁。

一维谐振子势中处于基态的粒子受到微扰
\[
H'(t) = \lambda x\, \delta(t)
\]
作用。求跃迁到各激发态的概率 $P_n$ ($n \geq 1$)、跃迁到所有激发态的总概率 $P_{\mathrm{exc}}$，以及仍处在基态的概率 $P_0$。

已知递推关系：
\[
x H_n(\xi) = \frac{1}{\sqrt{\alpha}} \left[ \sqrt{\frac{n}{2}} H_{n-1}(\xi) + \sqrt{\frac{n+1}{2}} H_{n+1}(\xi) \right], \qquad \alpha = \frac{m\omega}{\hbar}.
\]
\end{modelbox}

\subsection{考点快速分析}

\begin{tipbox}[考点快速分析]
\begin{itemize}
\item $\delta(t)$ 脉冲：时间演化算符 $U = \exp\!\bigl(-i\lambda x/\hbar\bigr)$，不是微扰展开
\item 位移算符：$x \propto a + a^\dagger$，指数化后正是相干态的生成算符
\item 相干态的 Fock 展开：$|\alpha\rangle = e^{-|\alpha|^2/2} \sum_n \dfrac{\alpha^n}{\sqrt{n!}} |n\rangle$
\item Poisson 分布：$P_n = \dfrac{e^{-\bar{n}} \bar{n}^n}{n!}$，平均量子数 $\bar{n} = \lambda^2/(2m\omega\hbar)$
\end{itemize}
\end{tipbox}

\subsection{$\delta$ 脉冲的时间演化算符}

\begin{derivationbox}[$\delta$ 脉冲的时间演化算符]
对 $H'(t) = \lambda x\, \delta(t)$，薛定谔方程在 $t=0$ 处积分：
\[
|\psi(0^+)\rangle - |\psi(0^-)\rangle = \frac{1}{i\hbar} \int_{-\epsilon}^{+\epsilon} \lambda x\, \delta(t') |\psi(t')\rangle dt' = \frac{\lambda x}{i\hbar} |\psi(0)\rangle.
\]
精确到 $O(\epsilon)$，得有限跳跃：
\[
|\psi(0^+)\rangle = e^{-i\lambda x/\hbar} |\psi(0^-)\rangle.
\]
因此时间演化算符为
\[
U = \exp\!\left( -\frac{i\lambda}{\hbar} x \right).
\]
\textbf{关键：}$\delta(t)$ 脉冲下算符是精确指数，不是微扰级数。这是与常微扰、简谐微扰的本质区别。
\end{derivationbox}

\subsection{位移算符与相干态}

\begin{derivationbox}[位移算符与相干态]
用升降算符表示位置：
\[
x = \sqrt{\frac{\hbar}{2m\omega}} \, (a + a^\dagger) \equiv x_0 (a + a^\dagger), \qquad x_0 = \sqrt{\frac{\hbar}{2m\omega}}.
\]
代入演化算符：
\[
U = \exp\!\left[ -\frac{i\lambda x_0}{\hbar} (a + a^\dagger) \right] = \exp\!\bigl( \alpha a^\dagger - \alpha^* a \bigr) \equiv D(\alpha),
\]
其中
\[
\alpha = -\frac{i\lambda x_0}{\hbar} = -\frac{i\lambda}{\hbar} \sqrt{\frac{\hbar}{2m\omega}}.
\]
$D(\alpha)$ 正是位移算符（displacement operator），作用在真空态上产生相干态：
\[
|\psi(0^+)\rangle = D(\alpha) |0\rangle = |\alpha\rangle.
\]

\subsubsection{位移算符的 BCH 处理}

利用 Baker--Campbell--Hausdorff 公式。由于 $[a, a^\dagger] = 1$，有
\[
e^{\alpha a^\dagger - \alpha^* a} = e^{-|\alpha|^2/2} \, e^{\alpha a^\dagger} e^{-\alpha^* a}.
\]
作用在 $|0\rangle$ 上，$e^{-\alpha^* a}|0\rangle = |0\rangle$（因为 $a|0\rangle = 0$），所以
\[
|\alpha\rangle = e^{-|\alpha|^2/2} e^{\alpha a^\dagger} |0\rangle = e^{-|\alpha|^2/2} \sum_{n=0}^\infty \frac{\alpha^n}{\sqrt{n!}} |n\rangle.
\]
\end{derivationbox}

\subsection{跃迁概率——Poisson 分布}

\begin{formulabox}[跃迁概率——Poisson 分布]
由相干态的 Fock 展开：
\[
\langle n|\alpha\rangle = e^{-|\alpha|^2/2} \frac{\alpha^n}{\sqrt{n!}}.
\]
跃迁概率：
\[
P_n = |\langle n|\alpha\rangle|^2 = \frac{e^{-|\alpha|^2} |\alpha|^{2n}}{n!}.
\]
令 $\bar{n} = |\alpha|^2 = \lambda^2 x_0^2/\hbar^2 = \dfrac{\lambda^2}{2m\omega\hbar}$，则
\[
P_n = \frac{e^{-\bar{n}} \, \bar{n}^n}{n!}, \qquad \bar{n} = \frac{\lambda^2}{2m\omega\hbar}.
\]
这正是 Poisson 分布！
\begin{itemize}
\item 基态概率：$P_0 = e^{-\bar{n}} = \exp\!\left( -\dfrac{\lambda^2}{2m\omega\hbar} \right)$
\item 激发态总概率：$P_{\mathrm{exc}} = \sum_{n=1}^\infty P_n = 1 - P_0 = 1 - e^{-\bar{n}}$
\item 平均激发数：$\langle n\rangle = \bar{n} = \lambda^2/(2m\omega\hbar)$
\item 方差：$\langle (\Delta n)^2\rangle = \bar{n}$（Poisson 分布的特征）
\end{itemize}
\end{formulabox}

\subsubsection{为什么出现 Poisson 分布？}

\begin{insightbox}[为什么出现 Poisson 分布？]
$\delta$ 脉冲给粒子一个瞬时的动量冲击 $\Delta p = \lambda$（由 $H' = \lambda x\delta(t)$，冲量 $\int F dt = -\partial H'/\partial x = \lambda\delta(t)$，所以动量跳跃为 $\lambda$）。

在相空间中，基态（最小不确定波包）被平移了一段距离，位置中心不动、动量中心跳跃。这个平移后的态正是相干态。

相干态是最经典的量子态，其数分布的涨落 $\Delta n = \sqrt{\bar{n}}$ 恰好是散粒噪声极限（shot-noise limit）。
\end{insightbox}

\subsubsection{用递推关系直接验证（不用相干态语言）}

\begin{derivationbox}[用递推关系直接验证]
对不熟悉相干态的考生，也可以直接用升降算符计算。由 $x = x_0(a + a^\dagger)$：
\[
U = e^{-i\lambda x/\hbar} = e^{\beta(a^\dagger - a)}, \qquad \beta = -\frac{i\lambda x_0}{\hbar}.
\]
利用 Glauber 公式 $e^{A+B} = e^A e^B e^{-[A,B]/2}$（当 $[A,B]$ 与 $A,B$ 都对易时），取 $A = \beta a^\dagger$，$B = -\beta a$，$[A,B] = -\beta^2[a^\dagger, a] = \beta^2$：
\[
U = e^{\beta^2/2} \, e^{\beta a^\dagger} e^{-\beta a}.
\]
作用在 $|0\rangle$ 上：$U|0\rangle = e^{\beta^2/2} e^{\beta a^\dagger} |0\rangle$。由于 $\beta$ 是纯虚数，$\beta^2 = -|\beta|^2$，所以 $e^{\beta^2/2} = e^{-|\beta|^2/2}$，与相干态结果一致。
\end{derivationbox}

\subsection{考场速记卡：三种脉冲形式的对比}

\begin{cheatbox}[考场速记卡：三种脉冲形式的对比]
\begin{center}
\begin{tabular}{ccc}
\hline
微扰形式 & 时间演化 & 结果特征 \\
\hline
$V\,\Theta(t)$（常微扰） & 微扰级数 & 费米黄金规则 \\
$V\,e^{-t/T}$（指数衰减） & 微扰级数 & Lorentz 线型 \\
$V\,\delta(t)$（瞬时脉冲） & 精确指数 $e^{-iV/\hbar}$ & 相干态/Poisson \\
\hline
\end{tabular}
\end{center}

\textbf{关键区分：}只有 $\delta(t)$ 能让 $U$ 写成闭式指数，因为积分 $\int \delta(t)\,dt$ 是有限的；其他形式都需要微扰展开。
\end{cheatbox}
